% --- IMPORTS ---
\documentclass[leqno]{article}
\usepackage{graphicx}
\usepackage{dsfont}
\usepackage{mathtools}
\usepackage{amssymb}
\usepackage{amsmath}
\usepackage{amsthm}
\usepackage{float}
\usepackage{ragged2e}
\usepackage{tensor}
\usepackage{fancyhdr}
\usepackage{empheq}
\usepackage[most]{tcolorbox}
\usepackage{enumitem}
\usepackage{dirtytalk}
\usepackage{marginnote}
\usepackage{microtype}
\usepackage[
  a4paper,
  left=0.8in,
  right=1in,
  top=1in,
  bottom=0.5in,
  headheight=14pt,
  headsep=0.4in
]{geometry}
\setlength{\parindent}{0cm}
% --- TikZ Packages ---
\usepackage{pgfplots}
\pgfplotsset{compat=1.18}
\usepgfplotslibrary{fillbetween, polar}
\usetikzlibrary{calc, positioning}
\usetikzlibrary{angles, quotes}
\usetikzlibrary{arrows.meta}
\usetikzlibrary{decorations.pathreplacing, decorations.markings}
\usetikzlibrary{patterns}
\usetikzlibrary{intersections}
% --- URL Packages ---
\usepackage[colorlinks=true,linkcolor=red,citecolor=magenta,urlcolor=black]{hyperref}%
\urlstyle{same}
% --- END OF IMPORTS ---

% --- CUSTOM COMMANDS ---
\RenewDocumentCommand{\marks}{o m}{%
  \IfValueTF{#1}%
    {\marginnote{\raggedright\textbf{[#2]}}[#1]}%
    {\marginnote{\raggedright\textbf{[#2]}}}%
}
\newcommand{\vspacerule}[2]{%
  \par\vspace*{#1}%
  \noindent\hrulefill\par
  \vspace*{#2}%
}
\definecolor{scarlet}{RGB}{205, 40, 75}
\newcommand{\questionitem}{%
  \item\phantomsection
  \addcontentsline{toc}{section}{Question \number\value{enumi}}%
}
\renewcommand{\contentsname}{Questions}
% --- END OF CUSTOM COMMANDS ---

% --- HEADER CONFIG ---
\pagestyle{fancy}
\fancyhead{}
\fancyfoot{}
\renewcommand{\headrulewidth}{0.9pt}
\lhead{Rational Functions}
\chead{}
\rhead{\href{https://danieljsmith.org}{danieljsmith.org}}
% --- END OF HEADER CONFIG ---

\begin{document}
\vspace*{20pt}
\tableofcontents
\newpage
\begin{enumerate}[
  label=\textbf{\arabic*.},
  leftmargin=*,
  topsep=0pt,
  partopsep=0pt,
  parsep=0pt
]

% ---- Question 1 ----
\questionitem
% [Source: _QBT__Rational_Functions, Q1]

A curve has equation
\[
y = 2 - \frac{3}{x+1}
\]
\begin{enumerate}[label=\textbf{(\alph*)}]
  \item State the equations of the horizontal and vertical asymptotes of this curve. \marks{2} \\
  \item The straight line $y=3x+c$ meets the curve. Show that the $x$-coordinates of the points of intersection satisfy
  \[
  3x^2+(c+1)x+(c+1)=0
  \]
  \marks[-10pt]{3}
  \item It is given that the line $y=3x+c$ is a tangent to the curve. \\
  \begin{enumerate}[label=(\roman*)]
    \item Without using calculus, find the two possible values of $c$.  \marks{3}
    \item Hence find the coordinates of each point of contact. \marks{4}
  \end{enumerate}
\end{enumerate}

\newpage

% ---- Question 2 ----
\questionitem
% [Source: _QBT__Rational_Functions, Q2]

The curve $C$ has equation $y = \dfrac{x+1}{x^2-4x+3}$. \\
\begin{enumerate}[label=\textbf{(\alph*)}]
  \item Find the equations of the asymptotes of $C$. \marks{2} \\
  \item Find the coordinates of any stationary points on $C$. \marks{4} \\
  \item Sketch $C$, stating the coordinates of the intersections with the axes. \marks{3} \\
  \item Sketch the curve with equation $y = \left|\dfrac{x+1}{x^2-4x+3}\right|$ and find in exact form the set of values of $x$ for which $2\lvert x+1\rvert > \lvert x^2-4x+3\rvert$. \marks{6} \\
\end{enumerate}

\newpage

% ---- Question 3 ----
\questionitem
% [Source: _QBT__Rational_Functions, Q3]

\begin{figure}[H]
    \centering
\begin{tikzpicture}[x=0.8cm,y=0.3cm]
\def\xmin{-7}
\def\xmax{8}
\def\ymin{-14}
\def\ymax{16}
\def\xa{-1}
\def\xb{2}
\begin{scope}
\clip (\xmin,\ymin) rectangle (\xmax,\ymax);
\draw[dashed, black] (\xmin,5) -- (\xmax,5);
\draw[dashed, black] (\xa,\ymin) -- (\xa,\ymax);
\draw[dashed, black] (\xb,\ymin) -- (\xb,\ymax);
\draw[thick, black, domain=-7:-1.05, samples=200, smooth] plot (\x, {(5*\x*\x - 2*\x + 1)/(\x*\x - \x - 2)});
\draw[thick, black, domain=-0.95:1.95, samples=300, smooth] plot (\x, {(5*\x*\x - 2*\x + 1)/(\x*\x - \x - 2)});
\draw[thick, black, domain=2.05:8, samples=200, smooth] plot (\x, {(5*\x*\x - 2*\x + 1)/(\x*\x - \x - 2)});
\end{scope}
\draw[-{Stealth[length=3mm]}, thick] (\xmin,0) -- (\xmax,0) node[right] {$x$};
\draw[-{Stealth[length=3mm]}, thick] (0,\ymin) -- (0,\ymax) node[above] {$y$};
\end{tikzpicture}
\end{figure}


The diagram shows the curve $C$ with equation
\[
y = \frac{5x^2 + ax - b}{x^2 + bx + a}
\]
where $a$ and $b$ are integers. \\

The vertical asymptotes of $C$ are $x=-1$ and $x=2$. \\
\begin{enumerate}[label=\textbf{(\alph*)}]
  \item State the equation of the asymptote of $C$ that is parallel to the $x$-axis. \marks{1} \\
  \item Use the vertical asymptotes to find the value of $a$ and the value of $b$. \marks{2} \\
  \item Hence, or otherwise, find the coordinates of the point where $C$ meets the $y$-axis. \marks{1} \\
  \item Without using calculus, show that the line $y=4$ does not intersect $C$. \marks{5} \\
\end{enumerate}

\newpage

% ---- Question 4 ----
\questionitem
% [Source: _QBT__Rational_Functions, Q4]

The function $f$ is defined by
\[
f(x)=\frac{-2x^2+11x-20}{2x-3} \qquad (x \in \mathbb{R},\ x \ne \tfrac{3}{2})
\]
\begin{enumerate}[label=\textbf{(\alph*)}]
  \item Find the range of $f$ \marks{5} \\
  \item Find the coordinates of the two stationary points of the graph of $y=f(x)$. \marks{2} \\
  \item Show that the graph of $y=f(x)$ has an oblique asymptote and find its equation. \marks{2} \\
  \item Sketch the graph of $y=f(x)$, showing clearly the asymptotes and stationary points. \marks{4} \\
\end{enumerate}

\newpage

% ---- Question 5 ----
\questionitem
% [Source: _QBT__Rational_Functions, Q5]

The curve $C$ has equation $y = \dfrac{x^2 + ax - 6}{x + 1}$, where $a > 0$. \\
\begin{enumerate}[label=\textbf{(\alph*)}]
  \item Find the equations of the asymptotes of $C$. \marks{3} \\
  \item Show that $C$ has no stationary points. \marks{4} \\
  \item Sketch $C$, stating the coordinates of the point of intersection with the $y$-axis and labelling the asymptotes. \marks{3} \\
  \item
  \begin{enumerate}[label=(\roman*)]
    \item Sketch the curve with equation $y = \left|\dfrac{x^2 + ax - 6}{x + 1}\right|$. \marks{2}
    \item On your sketch in part (i), draw the line $y = a + 1$. \marks{1}
    \item It is given that $\left|\dfrac{x^2 + ax - 6}{x + 1}\right| < a + 1$ for $-11 < x < -3$ and $0 < x < 4$. \\
    Find the value of $a$. \marks{2}
  \end{enumerate}
\end{enumerate}

\newpage

% ---- Question 6 ----
\questionitem
% [Source: _QBT__Rational_Functions, Q6]


The function $f$ is defined by
\[
f(x)=\frac{x^2+2x-3}{x^2+px+4} \qquad x \in \mathbb{R}
\]
where $p$ is a constant. \\

The graph of $y=f(x)$ has only one asymptote. \\
\begin{enumerate}[label=\textbf{(\alph*)}]
  \item Write down the equation of the asymptote. \marks{1} \\
  \item Find the set of possible values of $p$. \marks{4} \\
  \item Find the coordinates of the points at which the graph of $y=f(x)$ intersects the coordinate axes. \marks{3} \\

\begin{figure}[H]
    \centering
\begin{tikzpicture}
\pgfmathsetmacro{\xM}{4.5}
\pgfmathsetmacro{\yM}{4/3}
\pgfmathsetmacro{\kval}{0.2}
\begin{axis}[
    axis lines=center,
    xlabel={$x$},
    ylabel={$y$},
    xmin=-7, xmax=7,
    ymin=-1.2, ymax=1.6,
    xtick=\empty,
    ytick=\empty,
    clip=false,
    width=10cm,
    height=6.2cm,
    every axis x label/.style={at={(ticklabel* cs:1)}, anchor=west},
    every axis y label/.style={at={(ticklabel* cs:1)}, anchor=south}
]
\addplot[thick, black, samples=300, domain=-7:7] {(x^2 + 2*x - 3)/(x^2 - x + 4)};
\addplot[dashed, black, samples=2, domain=-7:7] {\kval};
\node[above right] at (axis cs:-7,0.6) {$C$};
\node[above right] at (axis cs:4.5,\kval) {$y=k$};
\node[below left] at (axis cs:0,0) {$O$};
\node[above] at (axis cs:\xM,\yM) {$M$};
\end{axis}
\end{tikzpicture}
\end{figure}

  
  \item A curve $C$ has equation
  \[
  y=\frac{x^2+2x-3}{x^2-x+4}
  \]
  The curve $C$ has a local maximum at the point $M$ as shown in the diagram. \\
  The line $y=k$ intersects curve $C$. \\
  \begin{enumerate}[label=(\roman*)]
    \item Show that
    \[
    15k^2-8k-16 \le 0
    \]
    \marks[-20pt]{5}
    \item Hence, find the $y$-coordinate of point $M$. \marks{2}
  \end{enumerate}
\end{enumerate}

\newpage

% ---- Question 7 ----
\questionitem
% [Source: _QBT__Rational_Functions, Q7]

A curve has equation
\[
y = \frac{x^2 - 9}{x^2 - 4}
\]
\begin{enumerate}[label=\textbf{(\alph*)}]
  \item Write down the equations of the asymptotes to the curve. \marks{3} \\
  \item Sketch the curve, indicating the coordinates of the points where the curve intersects the coordinate axes. \marks{5} \\
  \item Hence, or otherwise, solve the inequality
  \[
  0 < \frac{x^2 - 9}{x^2 - 4} < 1
  \]
  \marks[-20pt]{2}
\end{enumerate}

\newpage

% ---- Question 8 ----
\questionitem
% [Source: _QBT__Rational_Functions, Q8]

The curve $C$ has equation
\[
y=\frac{x^2+3}{x+1}
\]
\begin{enumerate}[label=\textbf{(\alph*)}]
  \item Find the equations of the asymptotes of $C$. \marks{3} \\
  \item Show that there is no point on $C$ for which $-6<y<2$. \marks{4} \\
  \item Sketch $C$. \marks{2} \\
  \item
  \begin{enumerate}[label=(\roman*)]
    \item Sketch the graphs of
    \[
    y=\left|\frac{x^2+3}{x+1}\right| \text{ and } y=|x-1|
    \]
    on the same axes, stating the coordinates of any intersections with the axes. \marks{4}
    \item Use your sketch to find the set of values of $c$ for which
    \[
    \left|\frac{x^2+3}{x+1}\right| \le |x-1|+c
    \]
    has no solution. \marks{1}
  \end{enumerate}
\end{enumerate}

\newpage

% ---- Question 9 ----
\questionitem
% [Source: _QBT__Rational_Functions, Q9]

Let $a$ be a positive constant. \\
\begin{enumerate}[label=\textbf{(\alph*)}]
  \item The curve $C_1$ has equation
  \[
  y=\frac{x^2-a^2}{x^2-4a^2}
  \]
  Sketch $C_1$. \marks{2} \\
  \end{enumerate}
  The curve $C_2$ has equation
  \[
  y=\left(\frac{x^2-a^2}{x^2-4a^2}\right)^2
  \]
  The curve $C_3$ has equation
  \[
  y=\left|\frac{x^2-a^2}{x^2-4a^2}\right|
  \]
  \\
  \begin{enumerate}[label=\textbf{(\alph*)}, resume]
  \item
  \begin{enumerate}[label=(\roman*)]
    \item Find the coordinates of all stationary points of $C_2$. \marks{3}
    \item Find also the coordinates of all points of intersection of $C_2$ and $C_3$. \marks{3} \\
  \end{enumerate}
  \item Sketch $C_2$ and $C_3$ on a single diagram, clearly identifying each curve. Hence find the set of values of $x$ for which
  \[
  \left(\frac{x^2-a^2}{x^2-4a^2}\right)^2 \leq \left|\frac{x^2-a^2}{x^2-4a^2}\right|
  \]
  \marks[-30pt]{5}
\end{enumerate}

\newpage

% ---- Question 10 ----
\questionitem
% [Source: _QBT__Rational_Functions, Q10]

A curve $C_1$ has equation
\[
y=-2+\frac{12}{x+3}
\]

\begin{enumerate}[label=\textbf{(\alph*)}]
  \item Write down the equations of the asymptotes of curve $C_1$. \marks{2} \\
  \item Sketch the graph of curve $C_1$. 
  Indicate the values of the intercepts of the curve with the axes. \marks{3}
  \item Hence, or otherwise, solve the inequality
  \[
  -2+\frac{12}{x+3}>0
  \]
  \marks[-20pt]{2}
  \item Curve $C_2$ is a reflection of curve $C_1$ in the line $y=-x$.
  Find an equation for curve $C_2$ in the form $y=f(x)$. \marks{3}
\end{enumerate}


\newpage

% ---- Question 11 ----
\questionitem
% [Source: _QBT__Rational_Functions, Q11]

The curve $C$ has equation
\[
y=\frac{x^2-6x+5}{x+2}
\] \\
\begin{enumerate}[label=\textbf{(\alph*)}]
  \item Find the equations of the asymptotes of $C$. \marks{3} \\
  \item Find the exact coordinates of the stationary points on $C$. \marks{4} \\
  \item Sketch $C$, stating the coordinates of any intersections with the axes. \marks{3} \\
  \item Sketch the curve with equation
  \[
  y=\left|\frac{x^2-6x+5}{x+2}\right|
  \] \\
  and find in exact form the set of values of $x$ for which
  \[
  \left|\frac{x^2-6x+5}{x+2}\right|<6
  \]
  \marks[-30pt]{5}
\end{enumerate}

\newpage

% ---- Question 12 ----
\questionitem
% [Source: _QBT__Rational_Functions, Q12]

A curve has equation $y=\dfrac{x-1}{x^2-x-6}$. \\
\begin{enumerate}[label=\textbf{(\alph*)}]
  \item
  \begin{enumerate}[label=(\roman*)]
    \item Write down the equations of the three asymptotes of the curve. \marks{3}
    \item Sketch the curve, showing the coordinates of any points of intersection with the coordinate axes. \marks[-10pt]{4} \\
  \end{enumerate}
  \item Hence, or otherwise, solve the inequality
  \[
  \frac{x-1}{x^2-x-6}>\frac{1}{2}
  \]
  \marks[-20pt]{3}
\end{enumerate}

\newpage

% ---- Question 13 ----
\questionitem
% [Source: _QBT__Rational_Functions, Q13]

The curve $C$ has equation
\[
y=\frac{(x-1)(x+3)}{x^2-2x+3}
\] \\
\begin{enumerate}[label=\textbf{(\alph*)}]
  \item Show that $C$ has no vertical asymptotes and state the equation of the horizontal asymptote of $C$. \marks{3} \\
  \item Find the coordinates of the stationary points on $C$. \marks{4} \\
  \item Sketch $C$, stating the coordinates of the intersections with the axes. \marks{3} \\
  \item Sketch the curve with equation
  \[
  y=\left|\frac{(x-1)(x+3)}{x^2-2x+3}\right|
  \] \\
  and state the set of values of $k$ for which
  \[
  \left|\frac{(x-1)(x+3)}{x^2-2x+3}\right|=k
  \]
  has $4$ distinct real solutions. \marks{2} \\
\end{enumerate}

\newpage

% ---- Question 14 ----
\questionitem
% [Source: _QBT__Rational_Functions, Q14]

The curve $C$ has equation
\[
y=\dfrac{-x^2+4x+5}{x+2}
\] \\
\begin{enumerate}[label=\textbf{(\alph*)}]
  \item Find the equations of the asymptotes of $C$. \marks{3} \\
  \item Find the coordinates of the stationary points of $C$. \marks{4} \\
  \item Sketch $C$, stating the coordinates of the intersections with the axes. \marks{3} \\
  \item Sketch the curve with equation
  \[
  y=\left|\dfrac{-x^2+4x+5}{x+2}\right|
  \]
  \marks[-20pt]{2}
  \item Using your sketch, or otherwise, find the set of values of $x$ for which
  \[
  \left|\dfrac{-x^2+4x+5}{x+2}\right|<2
  \]
  \marks[-30pt]{4}
\end{enumerate}
\end{enumerate}
\end{document}
